% Generated from 2/2-9.md; technical content preserved.
\exercise{۲.۹}
صورت \textbf{\lr{Exercise 2.9}} از فصل دوم کتاب درسی به بررسی کنترل ولتاژ آستانه از طریق اعمال بایاس به زیرلایه (\lr{Substrate Bias}) می‌پردازد:

\begin{quote}
\textbf{صورت مسئله:}

در برخی کاربردها، سابستریت (بدنه) به ولتاژی به نام \textbf{\lr{Substrate Bias}} متصل می‌شود تا ولتاژ آستانه ترانزیستورهای \lr{nMOS} تغییر داده شود. اگر هدف \textbf{افزایش ولتاژ آستانه (\(V_t\))} یک ترانزیستور \lr{nMOS} باشد، باید از \textbf{بایاس سابستریت مثبت} استفاده شود یا \textbf{منفی}؟

\end{quote}

\subsection*{۱. پاسخ مستقیم (\lr{Direct Answer})}
باید از \textbf{بایاس سابستریت منفی (\lr{Negative Substrate Bias})} یا اصطلاحاً \textbf{\lr{Reverse Body Bias} (\lr{RBB})} استفاده شود.

\subsection*{۲. اثبات و تحلیل بر پایه روابط تحلیلی (\lr{Mathematical Derivation})}

رابطه عمومی ولتاژ آستانه با در نظر گرفتن اثر بدنه (\lr{Body Effect}) به‌صورت زیر بیان می‌شود:
\[V_t = V_{t0} + \gamma \left( \sqrt{\phi_s + V_{SB}} - \sqrt{\phi_s} \right)\]

در این رابطه:
\begin{itemize}
\item \(V_{t0}\) ولتاژ آستانه نامی ترانزیستور در شرایط \(V_{SB} = 0\) است.
\item \(\gamma\) ضریب مثبت اثر بدنه (\lr{Body Effect Coefficient}) است.
\item \(\phi_s\) پتانسیل سطحی در شروع وارونگی قوی (مقداری مثبت) است.
\item \(V_{SB} = V_S - V_B\) ولتاژ سورس نسبت به بدنه (سابستریت) است.

\end{itemize}
در یک پیکربندی معمول، سورس ترانزیستور \lr{nMOS} به پتانسیل زمین مرجع متصل است (\(V_S = 0\text{\lr{ V}}\)). با اعمال ولتاژ بایاس سابستریت \(V_B\):
\[V_{SB} = 0 - V_B = -V_B\]

برای افزایش ولتاژ آستانه نسبت به حالت نامی (\(V_t > V_{t0}\))، لازم است ترم رادیکالی بزرگ‌تر از \(\sqrt{\phi_s}\) شود:
\[\sqrt{\phi_s + V_{SB}} > \sqrt{\phi_s} \implies V_{SB} > 0\]

با جایگذاری \(V_{SB} = -V_B\):
\[-V_B > 0 \implies \mathbf{V_B < 0}\]

بنابراین، اعمال یک ولتاژ منفی به سابستریت (\(V_B < 0\)) باعث مثبت شدن \(V_{SB}\) و در نتیجه افزایش ولتاژ آستانه \(V_t\) ترانزیستور \lr{nMOS} می‌شود.

\subsection*{۳. مکانیزم فیزیکی در ادوات نیمه‌هادی (\lr{Device Physics Intuition})}

\begin{itemize}
\item در ترانزیستور \lr{nMOS}، بدنه از جنس نیمه‌هادی نوع \lr{p} (\lr{p-type substrate}) است.
\item اتصال سابستریت به ولتاژ منفی، پیوند \lr{p-n} میان سابستریت و سورس/کانال را در \textbf{بایاس معکوس (\lr{Reverse Bias})} عمیق‌تری قرار می‌دهد.
\item با افزایش بایاس معکوس، حفره‌های اکثریت از ناحیه زیر گیت پس زده شده و \textbf{عرض ناحیه تخلیه (\lr{Depletion Region})} افزایش می‌یابد.
\item این امر بار یون‌های منفی غیرمتحرک اکسپتور (\lr{Acceptor Ions}) را در لایه تخلیه بیشتر می‌کند:
\[|Q_{dep}| = \sqrt{2 q \varepsilon_{si} N_A (\phi_s + V_{SB})}\]
\item برای ایجاد لایه وارونگی الکترونی (\lr{Inversion Layer}) در سطح کانال، بار مثبت روی الکترود گیت ابتدا باید میدان ناشی از این بارهای منفی لایه تخلیه را خنثی سازد.
\item چون بار منفی ناحیه تخلیه افزایش یافته است، به میدان الکتریکی عمودی و در نتیجه \textbf{ولتاژ گیت مثبت بزرگ‌تری} برای برقراری وارونگی قوی نیاز خواهد بود که این دقیقاً به معنای \textbf{افزایش ولتاژ آستانه (\(V_t\))} است.

\end{itemize}

\subsection*{۴. کاربرد در مدارهای مجتمع نانومتری (\lr{VLSI Context \& Trade-offs})}

این مکانیزم در تکنیک‌های مدیریت توان مدار مجتمع با عنوان \textbf{\lr{Adaptive Reverse Body Biasing} (\lr{RBB})} به کار می‌رود:

\begin{enumerate}
\item \textbf{کاهش توان مصرفی ایستا (\lr{Standby Leakage Reduction}):}
جریان نشتی زیر آستانه (\lr{Subthreshold Leakage}) به‌صورت نمایی با \(V_t\) رابطه عکس دارد:
\[I_{sub} \propto \exp\left(-\frac{V_t}{n v_T}\right)\]
در حالت آماده‌به‌کار (\lr{Standby/Sleep Mode})، با اعمال ولتاژ منفی به سابستریت (مثلاً \(V_B = -0.4\text{\lr{ V}}\))، ولتاژ آستانه بالا برده شده و جریان نشتی به شدت سرکوب می‌شود.

\item \textbf{محدودیت‌های بایاس معکوس:}
اگر ولتاژ منفی سابستریت بیش از حد بزرگ شود (مثلاً منفی‌تر از \(-1.2\text{\lr{ V}}\))، جریان‌های نشتی ناشی از شکست و تونل‌زنی پیوند نظیر \textbf{\lr{Band-to-Band Tunneling} (\lr{BTBT})} و \textbf{\lr{Gate-Induced Drain Leakage} (\lr{GIDL})} افزایش چشمگیری پیدا می‌کنند که اثر کاهش نشتی زیر آستانه را خنثی می‌سازد.

\end{enumerate}
